% --- Archon physics formalization source begin ---
% archon:physics
% archon:covers PhyXMiniProblems/problem_phyx_mini_0851.lean
% archon:source-report reports/phyx_mini/problem_phyx_mini_0851.source.json

\chapter{Physics problem phyx\_mini\_0851}
\label{ch:PhyXMiniProblems_problem_phyx_mini_0851}

\paragraph{Problem source.}
\#\# Question

Find the electric fluxe through surface 3.

\#\# Answer choices

- A: A: 80N m\textasciicircum{}2/C

- B: B: 10N m\textasciicircum{}2/C

- C: C: 20N m\textasciicircum{}2/C

- D: D: 0N m\textasciicircum{}2/C

\#\# Figure caption (auxiliary; use the image as primary evidence)

The image depicts a trapezoidal figure with dimensions and angles labeled. It appears to represent a physical or mathematical concept involving vectors and forces. Here's a detailed breakdown:

- The figure is a trapezoid with a base labeled as 4.0 m and a height of 2.0 m.
- The trapezoid is divided into four numbered sections (1, 2, 3, 4, and 5).
- A vector \textbackslash{}(\textbackslash{}vec\{E\}\textbackslash{}) is shown as a pink arrow pointing to the right, outside and parallel to the top line of the trapezoid. This vector is labeled \textbackslash{}(400 \textbackslash{}, \textbackslash{}text\{N/C\}\textbackslash{}).
- Another pink vector is drawn on the inside, along a diagonal line, also pointing to the right.
- A 30° angle is marked at the bottom right corner between two sides of the trapezoid.
- The interior sections are numbered with circles (from 1 to 5), potentially indicating different areas of interest or division.

\#\# Dataset metadata

Electrostatics; Spatial Relation Reasoning, Physical Model Grounding Reasoning

\paragraph{Recorded answer/context.}
D: D: 0N m\textasciicircum{}2/C

\paragraph{Figure/image path.}
/root/proposal\_for\_physic/science-mango/phyx\_mini\_run/phyx\_data/test\_image/851.png

\paragraph{Formalization target.}
create a compiling Lean file with sorry bodies at `PhyXMiniProblems/problem\_phyx\_mini\_0851.lean`. The Lean declarations must preserve the physical quantities, dimensions or dimensional roles, figure labels, governing-law hypotheses, and final relation expressed by this problem.
Use Mathlib/Physlib names found through LeanExplore where available. If a physics API is missing, introduce faithful local abstractions rather than scalar placeholder aliases.

\begin{theorem}[Physics formalization target]
\label{thm:physics:phyx_mini_0851:target}
\lean{PhyXMiniProblems.ProblemPhyXMini0851.problem_phyx_mini_0851}
\uses{def:physics:phyx-mini-0851:phyxminiproblems-problemphyxmini0851-electricfluxinnewtonsquaremeterspercoulomb, def:physics:phyx-mini-0851:phyxminiproblems-problemphyxmini0851-triangularprismfluxsetup, def:physics:phyx-mini-0851:phyxminiproblems-problemphyxmini0851-matchesproblemandprimaryfigure, def:physics:phyx-mini-0851:phyxminiproblems-problemphyxmini0851-satisfiestriangularprismgeometry, def:physics:phyx-mini-0851:phyxminiproblems-problemphyxmini0851-electricfieldalignedwithprismaxis, def:physics:phyx-mini-0851:phyxminiproblems-problemphyxmini0851-satisfiesuniformelectricfieldmodel, def:physics:phyx-mini-0851:phyxminiproblems-problemphyxmini0851-satisfiesplanarelectricfluxlaw, def:physics:phyx-mini-0851:phyxminiproblems-problemphyxmini0851-answerfluxinnewtonsquaremeterspercoulomb, def:physics:phyx-mini-0851:phyxminiproblems-problemphyxmini0851-recordeddatasetanswer}
The assigned autoformalize agent should translate this physics problem into Lean declarations in the covered file, with theorem and lemma proof bodies written as `by sorry`.
\end{theorem}
\begin{proof}
This is an autoformalization task, not a proof task. Produce statements that can later be proved by the physics prover without weakening the physical model.
\end{proof}

% --- Archon declaration topology begin ---
\section{Declaration topology}
\begin{definition}[electric Field Strength Dimension]
  \label{def:physics:phyx-mini-0851:phyxminiproblems-problemphyxmini0851-electricfieldstrengthdimension}
  \lean{PhyXMiniProblems.ProblemPhyXMini0851.electricFieldStrengthDimension}
  The physical dimension M L T\textasciicircum{}-2 C\textasciicircum{}-1 of electric-field strength
\end{definition}
\begin{proof}
  This is the declared model definition of electric Field Strength Dimension.
\end{proof}
\begin{definition}[area Dimension]
  \label{def:physics:phyx-mini-0851:phyxminiproblems-problemphyxmini0851-areadimension}
  \lean{PhyXMiniProblems.ProblemPhyXMini0851.areaDimension}
  The physical dimension L\textasciicircum{}2 of an oriented surface-area vector
\end{definition}
\begin{proof}
  This is the declared model definition of area Dimension.
\end{proof}
\begin{definition}[electric Flux Dimension]
  \label{def:physics:phyx-mini-0851:phyxminiproblems-problemphyxmini0851-electricfluxdimension}
  \lean{PhyXMiniProblems.ProblemPhyXMini0851.electricFluxDimension}
  \uses{def:physics:phyx-mini-0851:phyxminiproblems-problemphyxmini0851-electricfieldstrengthdimension, def:physics:phyx-mini-0851:phyxminiproblems-problemphyxmini0851-areadimension}
  The physical dimension M L\textasciicircum{}3 T\textasciicircum{}-2 C\textasciicircum{}-1 of electric flux
\end{definition}
\begin{proof}
  This is the declared model definition of electric Flux Dimension.
\end{proof}
\begin{definition}[Spatial Vector]
  \label{def:physics:phyx-mini-0851:phyxminiproblems-problemphyxmini0851-spatialvector}
  \lean{PhyXMiniProblems.ProblemPhyXMini0851.SpatialVector}
  Three-dimensional Cartesian vectors used for directions and SI components
\end{definition}
\begin{proof}
  This is the declared model definition of Spatial Vector.
\end{proof}
\begin{definition}[Length Quantity]
  \label{def:physics:phyx-mini-0851:phyxminiproblems-problemphyxmini0851-lengthquantity}
  \lean{PhyXMiniProblems.ProblemPhyXMini0851.LengthQuantity}
  A nonnegative, unit-independent physical length
\end{definition}
\begin{proof}
  This is the declared model definition of Length Quantity.
\end{proof}
\begin{definition}[Electric Field Strength Quantity]
  \label{def:physics:phyx-mini-0851:phyxminiproblems-problemphyxmini0851-electricfieldstrengthquantity}
  \lean{PhyXMiniProblems.ProblemPhyXMini0851.ElectricFieldStrengthQuantity}
  \uses{def:physics:phyx-mini-0851:phyxminiproblems-problemphyxmini0851-electricfieldstrengthdimension}
  A nonnegative, unit-independent electric-field strength
\end{definition}
\begin{proof}
  This is the declared model definition of Electric Field Strength Quantity.
\end{proof}
\begin{definition}[Electric Field Vector Quantity]
  \label{def:physics:phyx-mini-0851:phyxminiproblems-problemphyxmini0851-electricfieldvectorquantity}
  \lean{PhyXMiniProblems.ProblemPhyXMini0851.ElectricFieldVectorQuantity}
  \uses{def:physics:phyx-mini-0851:phyxminiproblems-problemphyxmini0851-electricfieldstrengthdimension, def:physics:phyx-mini-0851:phyxminiproblems-problemphyxmini0851-spatialvector}
  A unit-independent physical electric-field vector
\end{definition}
\begin{proof}
  This is the declared model definition of Electric Field Vector Quantity.
\end{proof}
\begin{definition}[Area Vector Quantity]
  \label{def:physics:phyx-mini-0851:phyxminiproblems-problemphyxmini0851-areavectorquantity}
  \lean{PhyXMiniProblems.ProblemPhyXMini0851.AreaVectorQuantity}
  \uses{def:physics:phyx-mini-0851:phyxminiproblems-problemphyxmini0851-areadimension, def:physics:phyx-mini-0851:phyxminiproblems-problemphyxmini0851-spatialvector}
  A unit-independent oriented physical area vector
\end{definition}
\begin{proof}
  This is the declared model definition of Area Vector Quantity.
\end{proof}
\begin{definition}[Signed Electric Flux Quantity]
  \label{def:physics:phyx-mini-0851:phyxminiproblems-problemphyxmini0851-signedelectricfluxquantity}
  \lean{PhyXMiniProblems.ProblemPhyXMini0851.SignedElectricFluxQuantity}
  \uses{def:physics:phyx-mini-0851:phyxminiproblems-problemphyxmini0851-electricfluxdimension}
  A signed, unit-independent electric flux
\end{definition}
\begin{proof}
  This is the declared model definition of Signed Electric Flux Quantity.
\end{proof}
\begin{definition}[length In Meters]
  \label{def:physics:phyx-mini-0851:phyxminiproblems-problemphyxmini0851-lengthinmeters}
  \lean{PhyXMiniProblems.ProblemPhyXMini0851.lengthInMeters}
  \uses{def:physics:phyx-mini-0851:phyxminiproblems-problemphyxmini0851-lengthquantity}
  Read a physical length in coherent-SI metres
\end{definition}
\begin{proof}
  This is the declared model definition of length In Meters.
\end{proof}
\begin{definition}[field Strength In Newtons Per Coulomb]
  \label{def:physics:phyx-mini-0851:phyxminiproblems-problemphyxmini0851-fieldstrengthinnewtonspercoulomb}
  \lean{PhyXMiniProblems.ProblemPhyXMini0851.fieldStrengthInNewtonsPerCoulomb}
  \uses{def:physics:phyx-mini-0851:phyxminiproblems-problemphyxmini0851-electricfieldstrengthquantity}
  Read an electric-field strength in newtons per coulomb
\end{definition}
\begin{proof}
  This is the declared model definition of field Strength In Newtons Per Coulomb.
\end{proof}
\begin{definition}[field Vector In Newtons Per Coulomb]
  \label{def:physics:phyx-mini-0851:phyxminiproblems-problemphyxmini0851-fieldvectorinnewtonspercoulomb}
  \lean{PhyXMiniProblems.ProblemPhyXMini0851.fieldVectorInNewtonsPerCoulomb}
  \uses{def:physics:phyx-mini-0851:phyxminiproblems-problemphyxmini0851-spatialvector, def:physics:phyx-mini-0851:phyxminiproblems-problemphyxmini0851-electricfieldvectorquantity}
  Read Cartesian electric-field components in newtons per coulomb
\end{definition}
\begin{proof}
  This is the declared model definition of field Vector In Newtons Per Coulomb.
\end{proof}
\begin{definition}[area Vector In Square Meters]
  \label{def:physics:phyx-mini-0851:phyxminiproblems-problemphyxmini0851-areavectorinsquaremeters}
  \lean{PhyXMiniProblems.ProblemPhyXMini0851.areaVectorInSquareMeters}
  \uses{def:physics:phyx-mini-0851:phyxminiproblems-problemphyxmini0851-spatialvector, def:physics:phyx-mini-0851:phyxminiproblems-problemphyxmini0851-areavectorquantity}
  Read an oriented area vector in square metres
\end{definition}
\begin{proof}
  This is the declared model definition of area Vector In Square Meters.
\end{proof}
\begin{definition}[electric Flux In Newton Square Meters Per Coulomb]
  \label{def:physics:phyx-mini-0851:phyxminiproblems-problemphyxmini0851-electricfluxinnewtonsquaremeterspercoulomb}
  \lean{PhyXMiniProblems.ProblemPhyXMini0851.electricFluxInNewtonSquareMetersPerCoulomb}
  \uses{def:physics:phyx-mini-0851:phyxminiproblems-problemphyxmini0851-signedelectricfluxquantity}
  Read signed electric flux in newton-square-metres per coulomb
\end{definition}
\begin{proof}
  This is the declared model definition of electric Flux In Newton Square Meters Per Coulomb.
\end{proof}
\begin{definition}[Surface Label]
  \label{def:physics:phyx-mini-0851:phyxminiproblems-problemphyxmini0851-surfacelabel}
  \lean{PhyXMiniProblems.ProblemPhyXMini0851.SurfaceLabel}
  The five numbered surfaces of the closed triangular prism
\end{definition}
\begin{proof}
  This is the declared model definition of Surface Label.
\end{proof}
\begin{definition}[number]
  \label{def:physics:phyx-mini-0851:phyxminiproblems-problemphyxmini0851-surfacelabel-number}
  \lean{PhyXMiniProblems.ProblemPhyXMini0851.SurfaceLabel.number}
  \uses{def:physics:phyx-mini-0851:phyxminiproblems-problemphyxmini0851-surfacelabel}
  The printed natural number associated with a surface label
\end{definition}
\begin{proof}
  This is the declared model definition of number.
\end{proof}
\begin{definition}[Prism Face Role]
  \label{def:physics:phyx-mini-0851:phyxminiproblems-problemphyxmini0851-prismfacerole}
  \lean{PhyXMiniProblems.ProblemPhyXMini0851.PrismFaceRole}
  Geometric role of a face in a triangular prism
\end{definition}
\begin{proof}
  This is the declared model definition of Prism Face Role.
\end{proof}
\begin{definition}[depicted Role]
  \label{def:physics:phyx-mini-0851:phyxminiproblems-problemphyxmini0851-surfacelabel-depictedrole}
  \lean{PhyXMiniProblems.ProblemPhyXMini0851.SurfaceLabel.depictedRole}
  \uses{def:physics:phyx-mini-0851:phyxminiproblems-problemphyxmini0851-surfacelabel, def:physics:phyx-mini-0851:phyxminiproblems-problemphyxmini0851-prismfacerole}
  The role of each face identified from the primary raster
\end{definition}
\begin{proof}
  This is the declared model definition of depicted Role.
\end{proof}
\begin{definition}[Triangular Prism Figure]
  \label{def:physics:phyx-mini-0851:phyxminiproblems-problemphyxmini0851-triangularprismfigure}
  \lean{PhyXMiniProblems.ProblemPhyXMini0851.TriangularPrismFigure}
  \uses{def:physics:phyx-mini-0851:phyxminiproblems-problemphyxmini0851-lengthquantity, def:physics:phyx-mini-0851:phyxminiproblems-problemphyxmini0851-electricfieldstrengthquantity, def:physics:phyx-mini-0851:phyxminiproblems-problemphyxmini0851-surfacelabel}
  Literal visual and dimensionful label content of image 851.png
\end{definition}
\begin{proof}
  This is the declared model definition of Triangular Prism Figure.
\end{proof}
\begin{definition}[Triangular Prism Flux Setup]
  \label{def:physics:phyx-mini-0851:phyxminiproblems-problemphyxmini0851-triangularprismfluxsetup}
  \lean{PhyXMiniProblems.ProblemPhyXMini0851.TriangularPrismFluxSetup}
  \uses{def:physics:phyx-mini-0851:phyxminiproblems-problemphyxmini0851-spatialvector, def:physics:phyx-mini-0851:phyxminiproblems-problemphyxmini0851-lengthquantity, def:physics:phyx-mini-0851:phyxminiproblems-problemphyxmini0851-electricfieldstrengthquantity, def:physics:phyx-mini-0851:phyxminiproblems-problemphyxmini0851-electricfieldvectorquantity, def:physics:phyx-mini-0851:phyxminiproblems-problemphyxmini0851-areavectorquantity, def:physics:phyx-mini-0851:phyxminiproblems-problemphyxmini0851-signedelectricfluxquantity, def:physics:phyx-mini-0851:phyxminiproblems-problemphyxmini0851-surfacelabel, def:physics:phyx-mini-0851:phyxminiproblems-problemphyxmini0851-triangularprismfigure}
  The physical apparatus and its independent observables. In particular, electricFluxThrough is not defined from an answer choice or from zero
\end{definition}
\begin{proof}
  This is the declared model definition of Triangular Prism Flux Setup.
\end{proof}
\begin{definition}[Matches Problem And Primary Figure]
  \label{def:physics:phyx-mini-0851:phyxminiproblems-problemphyxmini0851-matchesproblemandprimaryfigure}
  \lean{PhyXMiniProblems.ProblemPhyXMini0851.MatchesProblemAndPrimaryFigure}
  \uses{def:physics:phyx-mini-0851:phyxminiproblems-problemphyxmini0851-lengthinmeters, def:physics:phyx-mini-0851:phyxminiproblems-problemphyxmini0851-fieldstrengthinnewtonspercoulomb, def:physics:phyx-mini-0851:phyxminiproblems-problemphyxmini0851-triangularprismfluxsetup}
  All directly readable information from the problem and primary image. The angle is represented in radians, so the displayed 30° becomes pi / 6. No electric-flux value or answer choice occurs in these data
\end{definition}
\begin{proof}
  This is the declared model definition of Matches Problem And Primary Figure.
\end{proof}
\begin{definition}[Satisfies Triangular Prism Geometry]
  \label{def:physics:phyx-mini-0851:phyxminiproblems-problemphyxmini0851-satisfiestriangularprismgeometry}
  \lean{PhyXMiniProblems.ProblemPhyXMini0851.SatisfiesTriangularPrismGeometry}
  \uses{def:physics:phyx-mini-0851:phyxminiproblems-problemphyxmini0851-areavectorinsquaremeters, def:physics:phyx-mini-0851:phyxminiproblems-problemphyxmini0851-triangularprismfluxsetup}
  Coordinate-free triangular-prism geometry. A lateral face contains the prism axis, while each outward area vector spans a line orthogonal to its face tangent plane. This states geometry for every face rather than the requested flux value for face 3
\end{definition}
\begin{proof}
  This is the declared model definition of Satisfies Triangular Prism Geometry.
\end{proof}
\begin{definition}[Electric Field Aligned With Prism Axis]
  \label{def:physics:phyx-mini-0851:phyxminiproblems-problemphyxmini0851-electricfieldalignedwithprismaxis}
  \lean{PhyXMiniProblems.ProblemPhyXMini0851.ElectricFieldAlignedWithPrismAxis}
  \uses{def:physics:phyx-mini-0851:phyxminiproblems-problemphyxmini0851-fieldvectorinnewtonspercoulomb, def:physics:phyx-mini-0851:phyxminiproblems-problemphyxmini0851-triangularprismfluxsetup}
  The arrows in the raster show the electric field parallel to the prism axis and in its positive direction. The positive scalar is not fixed to a target flux and the field vector itself remains an independent physical quantity
\end{definition}
\begin{proof}
  This is the declared model definition of Electric Field Aligned With Prism Axis.
\end{proof}
\begin{definition}[Satisfies Uniform Electric Field Model]
  \label{def:physics:phyx-mini-0851:phyxminiproblems-problemphyxmini0851-satisfiesuniformelectricfieldmodel}
  \lean{PhyXMiniProblems.ProblemPhyXMini0851.SatisfiesUniformElectricFieldModel}
  \uses{def:physics:phyx-mini-0851:phyxminiproblems-problemphyxmini0851-fieldstrengthinnewtonspercoulomb, def:physics:phyx-mini-0851:phyxminiproblems-problemphyxmini0851-fieldvectorinnewtonspercoulomb, def:physics:phyx-mini-0851:phyxminiproblems-problemphyxmini0851-triangularprismfluxsetup}
  The unit-aware vector calibrates Physlib's electric field at every spacetime point, and its norm calibrates the independent physical strength
\end{definition}
\begin{proof}
  This is the declared model definition of Satisfies Uniform Electric Field Model.
\end{proof}
\begin{definition}[Satisfies Planar Electric Flux Law]
  \label{def:physics:phyx-mini-0851:phyxminiproblems-problemphyxmini0851-satisfiesplanarelectricfluxlaw}
  \lean{PhyXMiniProblems.ProblemPhyXMini0851.SatisfiesPlanarElectricFluxLaw}
  \uses{def:physics:phyx-mini-0851:phyxminiproblems-problemphyxmini0851-fieldvectorinnewtonspercoulomb, def:physics:phyx-mini-0851:phyxminiproblems-problemphyxmini0851-areavectorinsquaremeters, def:physics:phyx-mini-0851:phyxminiproblems-problemphyxmini0851-electricfluxinnewtonsquaremeterspercoulomb, def:physics:phyx-mini-0851:phyxminiproblems-problemphyxmini0851-triangularprismfluxsetup}
  The planar, uniform-field flux law Phi = E dot A, imposed on every face of the prism. This governing law does not single out surface 3 or assert zero
\end{definition}
\begin{proof}
  This is the declared model definition of Satisfies Planar Electric Flux Law.
\end{proof}
\begin{definition}[Answer Choice]
  \label{def:physics:phyx-mini-0851:phyxminiproblems-problemphyxmini0851-answerchoice}
  \lean{PhyXMiniProblems.ProblemPhyXMini0851.AnswerChoice}
  Labels of the four electric-flux choices in the source
\end{definition}
\begin{proof}
  This is the declared model definition of Answer Choice.
\end{proof}
\begin{definition}[answer Flux In Newton Square Meters Per Coulomb]
  \label{def:physics:phyx-mini-0851:phyxminiproblems-problemphyxmini0851-answerfluxinnewtonsquaremeterspercoulomb}
  \lean{PhyXMiniProblems.ProblemPhyXMini0851.answerFluxInNewtonSquareMetersPerCoulomb}
  \uses{def:physics:phyx-mini-0851:phyxminiproblems-problemphyxmini0851-answerchoice}
  Displayed flux in N m\textasciicircum{}2/C associated with each answer choice
\end{definition}
\begin{proof}
  This is the declared model definition of answer Flux In Newton Square Meters Per Coulomb.
\end{proof}
\begin{definition}[recorded Dataset Answer]
  \label{def:physics:phyx-mini-0851:phyxminiproblems-problemphyxmini0851-recordeddatasetanswer}
  \lean{PhyXMiniProblems.ProblemPhyXMini0851.recordedDatasetAnswer}
  \uses{def:physics:phyx-mini-0851:phyxminiproblems-problemphyxmini0851-answerchoice}
  The answer label recorded by the dataset
\end{definition}
\begin{proof}
  This is the declared model definition of recorded Dataset Answer.
\end{proof}
\begin{lemma}[field Vector mem surface Three Tangent Plane]
  \label{lem:physics:phyx-mini-0851:phyxminiproblems-problemphyxmini0851-fieldvector-mem-surfacethreetangentplane}
  \lean{PhyXMiniProblems.ProblemPhyXMini0851.fieldVector_mem_surfaceThreeTangentPlane}
  \uses{def:physics:phyx-mini-0851:phyxminiproblems-problemphyxmini0851-fieldvectorinnewtonspercoulomb, def:physics:phyx-mini-0851:phyxminiproblems-problemphyxmini0851-triangularprismfluxsetup, def:physics:phyx-mini-0851:phyxminiproblems-problemphyxmini0851-satisfiestriangularprismgeometry, def:physics:phyx-mini-0851:phyxminiproblems-problemphyxmini0851-electricfieldalignedwithprismaxis}
  Because surface 3 is lateral and the depicted field is axial, its field vector belongs to the surface-3 tangent plane. This is a geometric intermediate result and does not mention electric flux
\end{lemma}
\begin{proof}
  The conclusion follows from the governing relations and figure readouts named in the statement of field Vector mem surface Three Tangent Plane.
\end{proof}
\begin{lemma}[electric Flux Through Surface Three eq zero]
  \label{lem:physics:phyx-mini-0851:phyxminiproblems-problemphyxmini0851-electricfluxthroughsurfacethree-eq-zero}
  \lean{PhyXMiniProblems.ProblemPhyXMini0851.electricFluxThroughSurfaceThree_eq_zero}
  \uses{def:physics:phyx-mini-0851:phyxminiproblems-problemphyxmini0851-electricfluxinnewtonsquaremeterspercoulomb, def:physics:phyx-mini-0851:phyxminiproblems-problemphyxmini0851-triangularprismfluxsetup, def:physics:phyx-mini-0851:phyxminiproblems-problemphyxmini0851-satisfiestriangularprismgeometry, def:physics:phyx-mini-0851:phyxminiproblems-problemphyxmini0851-electricfieldalignedwithprismaxis, def:physics:phyx-mini-0851:phyxminiproblems-problemphyxmini0851-satisfiesplanarelectricfluxlaw}
  The tangent field is orthogonal to surface 3's outward area vector, so the planar flux law gives zero signed electric flux
\end{lemma}
\begin{proof}
  The conclusion follows from the governing relations and figure readouts named in the statement of electric Flux Through Surface Three eq zero.
\end{proof}
% --- Archon declaration topology end ---
% --- Archon physics formalization source end ---
